% Options for packages loaded elsewhere
\PassOptionsToPackage{unicode}{hyperref}
\PassOptionsToPackage{hyphens}{url}
\documentclass[
 12pt,
]{ctexart}
\usepackage{xcolor}
\usepackage{amsmath,amssymb}
\setcounter{secnumdepth}{-\maxdimen} % remove section numbering
\usepackage{iftex}
\ifPDFTeX
 \usepackage[T1]{fontenc}
 \usepackage[utf8]{inputenc}
 \usepackage{textcomp} % provide euro and other symbols
\else % if luatex or xetex
 \usepackage{unicode-math} % this also loads fontspec
 \defaultfontfeatures{Scale=MatchLowercase}
 \defaultfontfeatures[\rmfamily]{Ligatures=TeX,Scale=1}
\fi
\usepackage{lmodern}
\ifPDFTeX\else
 % xetex/luatex font selection
\fi
% Use upquote if available, for straight quotes in verbatim environments
\IfFileExists{upquote.sty}{\usepackage{upquote}}{}
\IfFileExists{microtype.sty}{% use microtype if available
 \usepackage[]{microtype}
 \UseMicrotypeSet[protrusion]{basicmath} % disable protrusion for tt fonts
}{}
\makeatletter
\@ifundefined{KOMAClassName}{% if non-KOMA class
 \IfFileExists{parskip.sty}{%
 \usepackage{parskip}
 }{% else
 \setlength{\parindent}{0pt}
 \setlength{\parskip}{6pt plus 2pt minus 1pt}}
}{% if KOMA class
 \KOMAoptions{parskip=half}}
\makeatother
\usepackage{longtable,booktabs,array}
\usepackage{calc} % for calculating minipage widths
% Correct order of tables after \paragraph or \subparagraph
\usepackage{etoolbox}
\makeatletter
\patchcmd\longtable{\par}{\if@noskipsec\mbox{}\fi\par}{}{}
\makeatother
% Allow footnotes in longtable head/foot
\IfFileExists{footnotehyper.sty}{\usepackage{footnotehyper}}{\usepackage{footnote}}
\makesavenoteenv{longtable}
\setlength{\emergencystretch}{3em} % prevent overfull lines
\providecommand{\tightlist}{%
 \setlength{\itemsep}{0pt}\setlength{\parskip}{0pt}}
\usepackage{bookmark}
\IfFileExists{xurl.sty}{\usepackage{xurl}}{} % add URL line breaks if available
\urlstyle{same}
\hypersetup{
 hidelinks,
 pdfcreator={LaTeX via pandoc}}

\author{SCX}
\date{2026}

\begin{document}

\section{Expert Governance Protocol ---
Formal Definition}\label{expert-governance-protocol-ux5f62ux5f0fux5316ux5b9aux4e49}

\begin{quote}
\textbf{SCX/EGP Pre-Distillation Expert Governance Protocol}\\
Version: v1.0\\
Last Updated:2026-06-26
\end{quote}

\begin{center}\rule{0.5\linewidth}{0.5pt}\end{center}

\subsection{Table of Contents}\label{ux76eeux5f55}

\begin{enumerate}
\def\labelenumi{\arabic{enumi}.}
\tightlist
\item
 \hyperref[1-ux7b26ux53f7ux4e0eux8bbeux5b9a]{Symboland设定}
\item
 \hyperref[2-ux4e94ux6b65-governance-protocol]{五步 Governance
 Protocol}

 \begin{itemize}
 \tightlist
 \item
 \hyperref[step-1-gauge-check-]{Step 1: Gauge Check}
 \item
 \hyperref[step-2-domain-certificate-]{Step 2: Domain Certificate}
 \item
 \hyperref[step-3-conflict-resolution-]{Step 3: Conflict Resolution}
 \item
 \hyperref[step-4-anchor-verification-]{Step 4: Anchor Verification}
 \item
 \hyperref[step-5-distillation-authorization-]{Step 5: Distillation
 Authorization}
 \end{itemize}
\item
 \hyperref[3-ux9608ux503cux9009ux62e9ux6307ux5f15]{threshold选择指引}
\item
 \hyperref[4-ux5931ux8d25ux56deux9000ux7b56ux7565]{失败Fallback strategy}
\item
 \hyperref[5-certification-theorem]{Certification Theorem}
\item
 \hyperref[6-aln-v3-ux56deux987eux6027ux9a8cux8bc1ux65b9ux6848]{AlN v3
 回顾性verification方案}
\item
 \hyperref[7-ux4e0e-scx-ux4ee3ux7801ux7684ux5b9eux73b0ux6620ux5c04]{and
 SCX 代码implementmap}
\end{enumerate}

\begin{center}\rule{0.5\linewidth}{0.5pt}\end{center}

\subsection{1. Symboland设定}\label{ux7b26ux53f7ux4e0eux8bbeux5b9a}

\subsubsection{1.1 Basic definitions}\label{ux57faux7840ux5b9aux4e49}

沿uses SCX frameworkdefinition(See \texttt{theory/README.md}):

\begin{longtable}[]{@{}
 >{\raggedright\arraybackslash}p{(\linewidth - 4\tabcolsep) * \real{0.3333}}
 >{\raggedright\arraybackslash}p{(\linewidth - 4\tabcolsep) * \real{0.3333}}
 >{\raggedright\arraybackslash}p{(\linewidth - 4\tabcolsep) * \real{0.3333}}@{}}
\toprule\noalign{}
\begin{minipage}[b]{\linewidth}\raggedright
Symbol
\end{minipage} & \begin{minipage}[b]{\linewidth}\raggedright
Meaning
\end{minipage} & \begin{minipage}[b]{\linewidth}\raggedright
Type
\end{minipage} \\
\midrule\noalign{}
\endhead
\bottomrule\noalign{}
\endlastfoot
\(\mathcal{X}\) & Input space(原子构型) & metricspace \\
\(\mathcal{Y}\) & Label space(energy/力) & \(\mathbb{R}^d\) \\
\(f^*: \mathcal{X} \to \mathcal{Y}\) & 真实Annotationfunction(DFT) & not yet知
oracle \\
\(f_m: \mathcal{X} \to \mathcal{Y}\) & 第 \(m\) 个 MLIP 专家 &
cancomputefunction,\(m \in [M]\) \\
\(M\) & 专家总数 & \(\mathbb{N}\) \\
\(\mathcal{S}\) & State space(we have限partition) & \(|\mathcal{S}| = K\) \\
\(s: \mathcal{X} \to \mathcal{S}\) & Statusmapfunction & can测function \\
\(s(x)\) & 构型 \(x\) 归属Status & \(\mathcal{S}\) \\
\(\ell: \mathcal{Y} \times \mathcal{Y} \to \mathbb{R}_{\ge 0}\) &
loss function(默认 MSE) & bounded:\(0 \le \ell \le L_{\max}\) \\
\end{longtable}

\subsubsection{1.2 Core quantities}\label{ux6838ux5fc3ux91cf}

\textbf{Statuscondition专家risk}(Definition 1, SCX):

\[R_m(s) = \mathbb{E}_{x \sim P(\cdot \mid s)}[\ell(f_m(x), f^*(x))]\]

\textbf{SCX can靠性}(Definition 2, SCX):

\[\text{SCX}_m(s) = P(\ell(f_m(x), f^*(x)) < \tau \mid x \in s)\]

\textbf{专家配auditingconflict}(Definition 3, ExpertConflict):

\[\text{Conflict}(m, n, s) = \mathbb{E}_{x \sim P(\cdot \mid s)}[|f_m(x) - f_n(x)|^2]\]

\subsubsection{1.3 Protocol
整体structure}\label{protocol-ux6574ux4f53ux7ed3ux6784}

Governance Protocol
isFromThe original专家system到蒸馏授权决策\textbf{linear判定流程}:

\[\mathcal{G}: \underbrace{(\{f_m\}, \mathcal{S}, \mathcal{D}_{\text{ref}})}_{\text{Input}} \xrightarrow{\text{Step 1-5}} \underbrace{\{\text{authorized}, \text{conditional}, \text{denied}\}}_{\text{Output}}\]

Input: - 专家集 \(\{f_m\}_{m=1}^M\)(can能异构:MACE, NEP, SevenNet,
\ldots) - State space \(\mathcal{S}\) 及Statusmap \(s(x)\) - 参考Data集
\(\mathcal{D}_{\text{ref}} = \{(x_i, y_i)\}_{i=1}^{N_{\text{ref}}}\)(DFT
label) - 蒸馏targetStatus \(s_{\text{target}} \in \mathcal{S}\)

Output: - \textbf{authorized}:allwe have检查via,can直接蒸馏 -
\textbf{conditional}:Partcondition不satisfybutcan补救(such as补anchorsTesting) -
\textbf{denied}:不can蒸馏,requires专家重新Trainingor更换专家

\begin{center}\rule{0.5\linewidth}{0.5pt}\end{center}

\subsection{2. 五步 Governance
Protocol}\label{ux4e94ux6b65-governance-protocol}

\subsubsection{Step 1: Gauge Check
(Gauge check)}\label{step-1-gauge-check-ux91cfux89c4ux68c0ux67e5}

\paragraph{Problem}\label{ux95eeux9898}

不同 MLIP 总energypredictexist insystem性energy偏移(gauge freedom).Step 1
detect并消除This些偏移,makingallwe have专家energypredictUnder统一参考scalingoncan比.

\paragraph{Formal Definition}\label{ux5f62ux5f0fux5316ux5b9aux4e49}

\textbf{Input}: - 专家集 \(\{f_m\}_{m=1}^M\) - 参考Status子集
\(\mathcal{S}_{\text{gauge}} \subseteq \mathcal{S}\)(选择energyalready知simpleStatus,such as体相晶格)
- 参考energy \(E_{\text{ref}}(s)\) auditingeach个
\(s \in \mathcal{S}_{\text{gauge}}\)(From DFT orExperiment获we get) - 容忍度
\(\epsilon_{\text{gauge}} > 0\)

\textbf{Gauge definition}:

each experts \(f_m\) energypredict分解provides:

\[f_m(x) = \underbrace{f_m^{\text{intrinsic}}(x)}_{\text{物理Part}} + \underbrace{\mu_m}_{\text{gauge offset}} + \underbrace{\eta_m(x)}_{\text{残差}}\]

where \(\mu_m\) isdepended on于专家systembias(gauge
shift).auditing力predict,\(\nabla_x \mu_m = 0\),Therefore力不受影响.but\textbf{总energy基准}Difference影响蒸馏学生auditingenergy学习.

\textbf{Gauge Check algorithm}:

\begin{enumerate}
\def\labelenumi{\arabic{enumi}.}
\item
 auditingeach个 \(s \in \mathcal{S}_{\text{gauge}}\),采样 \(N_{\text{gauge}}\)
 configurations \(\{x_{s,j}\}_{j=1}^{N_{\text{gauge}}}\)
\item
 computeeach expertsaveragingbias:

 \[\hat{\mu}_m(s) = \frac{1}{N_{\text{gauge}}} \sum_{j=1}^{N_{\text{gauge}}} \left[f_m(x_{s,j}) - E_{\text{ref}}(s)\right]\]
\item
 compute gauge consistent性metric:

 \[\text{GaugeVariance}(m) = \frac{1}{|\mathcal{S}_{\text{gauge}}|} \sum_{s \in \mathcal{S}_{\text{gauge}}} \left(\hat{\mu}_m(s) - \bar{\mu}_m\right)^2\]

 where
 \(\bar{\mu}_m = \frac{1}{|\mathcal{S}_{\text{gauge}}|} \sum_{s} \hat{\mu}_m(s)\)
\item
 判定condition:

 \[\text{gauge\_certified}(m) = \begin{cases}
 \text{true}, & \text{GaugeVariance}(m) \le \epsilon_{\text{gauge}} \\
 \text{false}, & \text{otherwise}
 \end{cases}\]

 (Low方差意味着偏移 \(\mu_m\) Undereach参考Status间consistent,is纯 gauge 偏移)
\end{enumerate}

\textbf{OutputType}:

\[\text{gauge\_certified}: [M] \to \{\text{true}, \text{false}\}\]

\textbf{Gauge Fixing(失败execute)}:

Ifexist in \(m\) makingwe get \(\text{gauge\_certified}(m) = \text{false}\),thenrequires
gauge fixing.definition校准function:

\[\tilde{f}_m(x) = f_m(x) - \frac{1}{|\mathcal{S}_{\text{gauge}}|} \sum_{s \in \mathcal{S}_{\text{gauge}}} \hat{\mu}_m(s)\]

i.e.,will专家 \(m\) energypredict减去itsauditing参考Statusaveragingsystembias.

校准after重新run Check.

\paragraph{andexisting工asrelationship}\label{ux4e0eux5df2ux6709ux5de5ux4f5cux7684ux5173ux7cfb}

\begin{itemize}
\tightlist
\item
 文献 ``Gauge dependence of total energies'' (arXiv:2504.05565, 2025)
 指出不同 DFT 泛函间exist in类似Problem,butnot yetprocess\textbf{跨 MLIP 架构} gauge
 auditing齐.
\item
 本Stepis SCX/EGP 原创contribution------Under蒸馏before强制多专家energyscaling统一.
\end{itemize}

\begin{center}\rule{0.5\linewidth}{0.5pt}\end{center}

\subsubsection{Step 2: Domain Certificate
(Domain certification)}\label{step-2-domain-certificate-ux57dfux8ba4ux8bc1}

\paragraph{Problem}\label{ux95eeux9898-1}

auditingtargetStatus
\(s_{\text{target}}\),evaluateeach expertsUndertheStatusoncan靠性.不can靠专家不should参and蒸馏.

\paragraph{Formal Definition}\label{ux5f62ux5f0fux5316ux5b9aux4e49-1}

\textbf{Input}: - 专家集 \(\{f_m\}_{m=1}^M\) - targetStatus
\(s_{\text{target}} \in \mathcal{S}\) - StatusAnnotationData
\(\mathcal{D}_{s} = \{(x_i, y_i) \mid s(x_i) = s_{\text{target}}\}\)(can能we haveor没we have)
- Lowcan靠性threshold \(\theta_{\text{low}} \in (0, 1)\) - Highcan靠性threshold
\(\theta_{\text{high}} \in (0, 1)\)(\(\theta_{\text{high}} > \theta_{\text{low}}\))
- Lossthreshold \(\tau > 0\) - 最小sample数 \(n_{\min} \in \mathbb{N}\)

\textbf{Reliability estimation}:

auditingeach experts \(m\),Under \(s_{\text{target}}\) onestimate SCX can靠性:

\[\widehat{\text{SCX}}_m(s_{\text{target}}) = \begin{cases}
\displaystyle \frac{1}{|\mathcal{D}_{s}|} \sum_{(x, y) \in \mathcal{D}_{s}} \mathbf{1}\{\ell(f_m(x), y) < \tau\}, & |\mathcal{D}_{s}| \ge n_{\min} \\[1.5em]
\text{Bayesian shrinkage estimate}, & |\mathcal{D}_{s}| < n_{\min}
\end{cases}\]

小samplemakinguses James-Stein Contractionestimate(auditingshould
\texttt{ExpertReliability.estimate\_small\_sample}):

\[\widehat{\text{SCX}}_m^{\text{JS}}(s) = \lambda \cdot \widehat{\text{SCX}}_m(s) + (1 - \lambda) \cdot \overline{\text{SCX}}_m\]

where
\(\lambda = \frac{|\mathcal{D}_{s}|}{|\mathcal{D}_{s}| + n_{\min}}\),\(\overline{\text{SCX}}_m\)
is专家 \(m\) Underallwe haveStatuson全局averagingcan靠性.

\textbf{Decision logic}:

\[\text{domain\_certified}(m, s) = \begin{cases}
\text{true}, & \widehat{\text{SCX}}_m(s) \ge \theta_{\text{high}} \\[0.5em]
\text{conditional}, & \theta_{\text{low}} \le \widehat{\text{SCX}}_m(s) < \theta_{\text{high}} \\[0.5em]
\text{false}, & \widehat{\text{SCX}}_m(s) < \theta_{\text{low}}
\end{cases}\]

\begin{itemize}
\tightlist
\item
 \textbf{true}: 专家UndertargetStatusonHigh度can靠,can参and蒸馏
\item
 \textbf{conditional}: can靠性不足以Independentcontribution,butcan参and多专家集成
\item
 \textbf{false}: 专家UndertheStatus不can信,exclude出蒸馏
\end{itemize}

\textbf{OutputType}:

\[\text{domain\_certified}: [M] \times \mathcal{S} \to \{\text{true}, \text{false}, \text{conditional}\}\]

\paragraph{and SCX
代码auditingshould}\label{ux4e0e-scx-ux4ee3ux7801ux7684ux5bf9ux5e94}

Step 2 直接调uses \texttt{ExpertReliability.estimate()} we get到
\texttt{SCX\_matrix},where \texttt{SCX\_matrix{[}m,\ k{]}} i.e.,
\(\widehat{\text{SCX}}_m(s_k)\).

\begin{center}\rule{0.5\linewidth}{0.5pt}\end{center}

\subsubsection{Step 3: Conflict Resolution
(conflictResolution)}\label{step-3-conflict-resolution-ux51b2ux7a81ux89e3ux51b3}

\paragraph{Problem}\label{ux95eeux9898-2}

auditingtargetStatus \(s_{\text{target}}\),检查via Step 2
认证专家之间whetherexist in显著predictconflict.conflict区域requires仲裁.

\paragraph{Formal Definition}\label{ux5f62ux5f0fux5316ux5b9aux4e49-2}

\textbf{Input}: - already认证专家子集
\(\mathcal{E}_{\text{cert}} \subseteq \{f_1, \dots, f_M\}\)(Step 2 Output
\(\text{domain\_certified} \in \{\text{true}, \text{conditional}\}\)) -
targetStatus \(s_{\text{target}}\) - Statuswithinnot yetAnnotation构型
\(X_s = \{x \mid s(x) = s_{\text{target}}, x \notin \mathcal{D}_{\text{ref}}\}\)
- conflictthreshold \(\theta_{\text{conflict}} > 0\)

\textbf{Conflict metric}:

auditingeach一auditing认证专家
\((m, n) \in \mathcal{E}_{\text{cert}} \times \mathcal{E}_{\text{cert}}\),compute:

\[\text{Conflict}(m, n, s) = \frac{1}{|X_s|} \sum_{x \in X_s} |f_m(x) - f_n(x)|^2\]

\textbf{Conflict determination}:

\[\text{HasConflict}(s) = \begin{cases}
\text{true}, & \exists (m, n): \text{Conflict}(m, n, s) > \theta_{\text{conflict}} \\
\text{false}, & \text{otherwise}
\end{cases}\]

If \(\text{HasConflict}(s) = \text{false}\),then
\(\text{conflict\_resolved} = \text{true}\)(NoneconflictrequiresResolution).

If \(\text{HasConflict}(s) = \text{true}\),requires进一步AnalysisconflictType:

\textbf{conflictType诊断}:

\begin{longtable}[]{@{}
 >{\raggedright\arraybackslash}p{(\linewidth - 4\tabcolsep) * \real{0.3333}}
 >{\raggedright\arraybackslash}p{(\linewidth - 4\tabcolsep) * \real{0.3333}}
 >{\raggedright\arraybackslash}p{(\linewidth - 4\tabcolsep) * \real{0.3333}}@{}}
\toprule\noalign{}
\begin{minipage}[b]{\linewidth}\raggedright
Type
\end{minipage} & \begin{minipage}[b]{\linewidth}\raggedright
condition
\end{minipage} & \begin{minipage}[b]{\linewidth}\raggedright
Meaning
\end{minipage} \\
\midrule\noalign{}
\endhead
\bottomrule\noalign{}
\endlastfoot
\textbf{Gauge Residual gauge conflict} &
\(\exists m,n: \text{Conflict}(m,n,s) > \theta_{\text{conflict}}\)
butenergy差approximateconstant & Step 1 not yetcompletely消除 gauge 残差 \\
\textbf{Structure-sensitive conflict} & conflictonly发生Under特定子区域 \(s' \subset s\) &
专家auditing局部structureType选择性失效 \\
\textbf{全Statusconflict} & conflictUndertheStatusallwe havesampleonallexist in & at least一 expertsUnder
\(s\) completely不can靠 \\
\end{longtable}

\textbf{Decision logic}:

\[\text{conflict\_resolved} = \begin{cases}
\text{true}, & \text{HasConflict}(s) = \text{false} \\[0.5em]
\text{true}, & \text{conflict经仲裁aftercanResolution(seeunder方)} \\[0.5em]
\text{needs\_anchor}, & \text{requires DFT anchor 仲裁}
\end{cases}\]

\textbf{Arbitration method}(conflictcanResolution):

IfconflictTypeprovides ``Gauge Residual gauge conflict'',返回 Step 1 重新校准. IfconflictTypeprovides
``Structure-sensitive conflict'',willconflict子区域 \(s'\) 标记providesrequires DFT anchor(Step
4)并excludeit,Under剩余区域on \(\text{conflict\_resolved} = \text{true}\).
IfconflictTypeprovides
``全Statusconflict'',\(\text{conflict\_resolved} = \text{needs\_anchor}\).

\textbf{OutputType}:

\[\text{conflict\_resolved}: \mathcal{S} \to \{\text{true}, \text{false}, \text{needs\_anchor}\}\]

\paragraph{and SCX
代码auditingshould}\label{ux4e0e-scx-ux4ee3ux7801ux7684ux5bf9ux5e94-1}

Step 3 makinguses \texttt{ExpertConflict.conflict\_matrix()} computeconflicttensor
\(D \in \mathbb{R}^{M \times M \times K}\),uses
\texttt{ExpertConflict.conflict\_score()} acquire归一化conflict评分,uses
\texttt{ExpertConflict.arbitrate()} UnderconflictStatusunder选择仲裁策略.

\begin{center}\rule{0.5\linewidth}{0.5pt}\end{center}

\subsubsection{Step 4: Anchor Verification
(anchorsverification)}\label{step-4-anchor-verification-ux951aux70b9ux9a8cux8bc1}

\paragraph{Problem}\label{ux95eeux9898-3}

auditing Step 3 标记provides \(\text{needs\_anchor}\) conflictStatus,uses DFT
computeanchors(anchor frames)asprovides ground truth 仲裁.

\paragraph{Formal Definition}\label{ux5f62ux5f0fux5316ux5b9aux4e49-3}

\textbf{Input}: - conflictStatus \(s_{\text{conflict}}\) 列表(Step 3 Output
\(\text{needs\_anchor}\) Status) - anchors预算
\(N_{\text{anchor}} \in \mathbb{N}\)(允许总 DFT compute数) -
anchorsPrecisionthreshold \(\theta_{\text{anchor}} > 0\) - anchors选择策略
\(\pi_{\text{anchor}}\)

\textbf{anchors选择}:

auditingeach个conflictStatus \(s\),allocateanchors数:

\[n_{\text{anchor}}(s) = \left\lceil N_{\text{anchor}} \cdot \frac{\rho(s) \cdot \text{ConflictScore}(s)}{\sum_{s'} \rho(s') \cdot \text{ConflictScore}(s')} \right\rceil\]

where \(\rho(s) = P(x \in s)\) isStatusprobabilistic,\(\text{ConflictScore}(s)\) is
Step 3 conflict度归一化metric.

anchors选择策略 \(\pi_{\text{anchor}}\)------in state \(s\) within选取
\(n_{\text{anchor}}(s)\) configurations \(x_{\text{anchor}}\) Method:

\[\pi_{\text{anchor}}: \begin{cases}
\text{is maximized分歧采样}: & \displaystyle \arg\max_{x \in s} \text{Var}_{m \in \mathcal{E}_{\text{cert}}}[f_m(x)] \\[0.8em]
\text{FPS(最远点采样)}: & \displaystyle \arg\max_{x \in s} \min_{x' \in \mathcal{D}_{\text{anchor}}} d(x, x')
\end{cases}\]

推荐makinguses\textbf{is maximized分歧采样}------选择专家predict方差is maximized构型,This些构型最能区分不同专家表现.

\textbf{anchorsverification}:

auditingeach个all选anchors \(x_{\text{anchor}}\),compute DFT label
\(y_{\text{anchor}} = f^*(x_{\text{anchor}})\).

auditingeach experts \(m\),verificationitsUnderanchorson误差:

\[\text{AnchorMAE}(m, s) = \frac{1}{n_{\text{anchor}}(s)} \sum_{j=1}^{n_{\text{anchor}}(s)} \left| f_m(x_{\text{anchor},j}) - y_{\text{anchor},j} \right|\]

全局anchorsverification:

\[\text{anchor\_verified} = \begin{cases}
\text{true}, & \max_{m \in \mathcal{E}_{\text{cert}}} \text{AnchorMAE}(m, s) \le \theta_{\text{anchor}} \text{ for all conflicting } s \\[0.5em]
\text{false}, & \text{otherwise}
\end{cases}\]

If \(\text{anchor\_verified} = \text{false}\),identify出表现最差专家
\(m_{\text{fail}} = \arg\max_m \text{AnchorMAE}(m, s)\),willitsexclude出蒸馏.

\textbf{校正after仲裁规then}(Ifanchorsverificationviabut仍we haveconflict):

Underanchors约束under,auditing专家weights进行校准:

\[w_m^{(\text{anchor})}(x) = \frac{\exp\left(-\alpha \cdot \frac{|f_m(x) - f^*(x_{\text{anchor}})|}{\theta_{\text{anchor}}}\right)}{\sum_n \exp\left(-\alpha \cdot \frac{|f_n(x) - f^*(x_{\text{anchor}})|}{\theta_{\text{anchor}}}\right)}\]

This确保了anchors附近predict以 DFT provides基准加权.

\textbf{OutputType}:

\[\text{anchor\_verified}: \mathcal{P}(\mathcal{S}) \to \{\text{true}, \text{false}\}\]

\begin{center}\rule{0.5\linewidth}{0.5pt}\end{center}

\subsubsection{Step 5: Distillation Authorization
(蒸馏授权)}\label{step-5-distillation-authorization-ux84b8ux998fux6388ux6743}

\paragraph{Problem}\label{ux95eeux9898-4}

聚合before 4 步Output,做出最终蒸馏授权决策.

\paragraph{Formal Definition}\label{ux5f62ux5f0fux5316ux5b9aux4e49-4}

\textbf{Input}: -
\(\text{gauge\_certified}: [M] \to \{\text{true}, \text{false}\}\)(Step
1) -
\(\text{domain\_certified}: [M] \times \mathcal{S} \to \{\text{true}, \text{false}, \text{conditional}\}\)(Step
2) -
\(\text{conflict\_resolved}: \mathcal{S} \to \{\text{true}, \text{false}, \text{needs\_anchor}\}\)(Step
3) -
\(\text{anchor\_verified}: \mathcal{P}(\mathcal{S}) \to \{\text{true}, \text{false}\}\)(Step
4) - targetStatus \(s_{\text{target}}\)

\textbf{授权判定}:

首先definition Step 2 聚合Domain certification:

\[\text{DomainCertAgg}(s) = \begin{cases}
\text{true}, & \exists m: \text{domain\_certified}(m, s) = \text{true} \\[0.5em]
\text{false}, & \nexists m: \text{domain\_certified}(m, s) \in \{\text{true}, \text{conditional}\}
\end{cases}\]

i.e.,at leastwe have一 expertsUnder \(s\) onis bounded below by认证providescan靠orusescondition参and.

\[\text{authorized} = \underbrace{\left(\forall m: \text{gauge\_certified}(m) = \text{true}\right)}_{\text{allwe have专家 gauge alreadyauditing齐}} \land \underbrace{\text{DomainCertAgg}(s_{\text{target}})}_{\text{targetStatuswe havecan靠专家}} \land \underbrace{\text{conflict\_resolved}(s_{\text{target}}) \in \{\text{true}, \text{needs\_anchor}\}}_{\text{conflictalreadyResolutionorusesanchorscanResolution}} \land \underbrace{\text{anchor\_verified}(\{s_{\text{target}}\})}_{\text{anchorsverificationvia}}\]

\textbf{OutputType}:

\[\text{authorized}: \{\text{true}, \text{false}\}\]

\textbf{Special cases}:

\begin{itemize}
\tightlist
\item
 Ifallwe have专家
 \(\text{gauge\_certified} = \text{false}\):\textbf{denied}------must先execute
 gauge fixing
\item
 If \(\text{conflict\_resolved} = \text{needs\_anchor}\)
 butanchors预算不足:返回 \textbf{conditional} with anchors建议
\item
 If \(\text{anchor\_verified} = \text{false}\)
 butexclude失败专家after重新viaallwe have检查:返回 \textbf{authorized} with
 exclude列表
\end{itemize}

\textbf{Distillation authorization certificate}(Protocol 标准Output):

\begin{verbatim}
DistillationAuthorization {
 status: authorized | conditional | denied,
 expert_subset: list[int],  // 参and蒸馏专家 ID
 excluded_experts: list[int],  // is bounded below byexclude专家及原because
 anchor_frames: int,  // makingusesanchors帧数
 gauge_offsets: dict[float],  // shoulduses gauge 校正量
 certificate_bound: float,  // ε(N_anchor, δ) 保证界
 confidence: float,  // confidence 1-δ
}
\end{verbatim}

\begin{center}\rule{0.5\linewidth}{0.5pt}\end{center}

\subsection{3. threshold选择指引}\label{ux9608ux503cux9009ux62e9ux6307ux5f15}

\subsubsection{3.1 threshold汇总}\label{ux9608ux503cux6c47ux603b}

\begin{longtable}[]{@{}
 >{\raggedright\arraybackslash}p{(\linewidth - 6\tabcolsep) * \real{0.2000}}
 >{\raggedright\arraybackslash}p{(\linewidth - 6\tabcolsep) * \real{0.2000}}
 >{\raggedright\arraybackslash}p{(\linewidth - 6\tabcolsep) * \real{0.2667}}
 >{\raggedright\arraybackslash}p{(\linewidth - 6\tabcolsep) * \real{0.3333}}@{}}
\toprule\noalign{}
\begin{minipage}[b]{\linewidth}\raggedright
threshold
\end{minipage} & \begin{minipage}[b]{\linewidth}\raggedright
Symbol
\end{minipage} & \begin{minipage}[b]{\linewidth}\raggedright
Recommended value
\end{minipage} & \begin{minipage}[b]{\linewidth}\raggedright
Selection rationale
\end{minipage} \\
\midrule\noalign{}
\endhead
\bottomrule\noalign{}
\endlastfoot
Lossthreshold & \(\tau\) & \(1\text{ meV/atom}\)(energy)or
\(0.05\text{ eV/angstrom}\)(力) & 物理onwe have意义Precisionboundary;DFT
itself数值Precision约 \(0.1\text{ meV/atom}\) \\
Lowcan靠性threshold & \(\theta_{\text{low}}\) & \(0.7\) &
Low于this值专家UndertargetStatuson误差probabilistic \textgreater{}
30\%,不适合Independent参and蒸馏 \\
Highcan靠性threshold & \(\theta_{\text{high}}\) & \(0.9\) &
High于this值专家canIndependentcontributionpredict \\
conflictthreshold & \(\theta_{\text{conflict}}\) &
\(5\text{ meV/atom}\)(byenergyMSE) & 两倍于典型 MLIP 
MAE(\textasciitilde2-3
meV/atom),conflict超过this值意味着专家间exist in实质性分歧 \\
anchorsPrecisionthreshold & \(\theta_{\text{anchor}}\) & \(2\text{ meV/atom}\)(MAE)
& anchorsverification允许is maximized专家-DFT bias,this值should略Low于 \(\tau\) \\
Gauge 方差threshold & \(\epsilon_{\text{gauge}}\) & \(0.5\text{ meV/atom}\) &
参考Status间 gauge 偏移方差on限,过大方差意味着非consistent偏移(非纯
gauge) \\
最小sample数 & \(n_{\min}\) & \(5\) &
Statuswithinsample数Low于this值requiresmakinguses贝叶斯Contractionestimate \\
\end{longtable}

\subsubsection{3.2
threshold选择Physical basis}\label{ux9608ux503cux9009ux62e9ux7684ux7269ux7406ux4f9dux636e}

Under MLIP scenario,energyand力Precisionthresholdshouldthat the optimal selection based on以under考量:

\begin{enumerate}
\def\labelenumi{\arabic{enumi}.}
\item
 \textbf{DFT 数值Precision}: 典型 DFT compute数值convergentPrecision约provides
 \(0.1\text{ meV/atom}\)(energy)and
 \(0.01\text{ eV/angstrom}\)(力).\(\tau\)
 should设于this值or more以避免drift in''虚假''专家Error.
\item
 \textbf{MLIP 典型Precisionscope}: asbefore MLIP auditingenergy MAE typicallyUnder
 \(1\text{--}20\text{ meV/atom}\) scope.\(\theta_{\text{conflict}}\)
 设providesthisscope值(\(5\text{ meV/atom}\))can合理区分正常分歧and异常conflict.
\item
 \textbf{力andenergy权衡}:
 力thresholdtypicallyshould比energythreshold宽松,becauseprovides力数值noisetypically更大.建议
 \(\tau_{\text{force}} = 10 \times \tau_{\text{energy}}\).
\item
 \textbf{自适shouldthreshold}:
 Under实际操as,can以that the optimal selection based onverification集on专家performancedistribution自适should设定threshold:
 \[\theta_{\text{low}} = \text{median}\{\widehat{\text{SCX}}_m(s)\} - 2 \cdot \text{MAD}\{\widehat{\text{SCX}}_m(s)\}\]
 where MAD provides位数绝auditingbias.
\end{enumerate}

\begin{center}\rule{0.5\linewidth}{0.5pt}\end{center}

\subsection{4. 失败Fallback strategy}\label{ux5931ux8d25ux56deux9000ux7b56ux7565}

\begin{longtable}[]{@{}
 >{\raggedright\arraybackslash}p{(\linewidth - 4\tabcolsep) * \real{0.2308}}
 >{\raggedright\arraybackslash}p{(\linewidth - 4\tabcolsep) * \real{0.3846}}
 >{\raggedright\arraybackslash}p{(\linewidth - 4\tabcolsep) * \real{0.3846}}@{}}
\toprule\noalign{}
\begin{minipage}[b]{\linewidth}\raggedright
Step
\end{minipage} & \begin{minipage}[b]{\linewidth}\raggedright
Failure mode
\end{minipage} & \begin{minipage}[b]{\linewidth}\raggedright
Fallback strategy
\end{minipage} \\
\midrule\noalign{}
\endhead
\bottomrule\noalign{}
\endlastfoot
\textbf{Step 1} & \(\text{gauge\_certified}(m) = \text{false}\) & execute
gauge fixing:\(\tilde{f}_m = f_m - \bar{\mu}_m\);重检 \\
\textbf{Step 1} & 参考Status \(\mathcal{S}_{\text{gauge}}\) NoneData &
From任意Status采样 \(N_{\text{gauge}}\) configurations,uses DFT compute参考energy \\
\textbf{Step 2} & \(\text{domain\_certified}(m, s) = \text{false}\)
auditingallwe have \(m\) & expandStatusscope(merge相邻Status以increasesample),or降Low
\(\theta_{\text{low}}\).If仍失败:requiresprovides新StatusSupplement DFT label \\
\textbf{Step 2} & \(|\mathcal{D}_{s}| = 0\)(全新Status) & makinguses Bayesian
shrinkage 到全局all值;标记provides \(\text{conditional}\) \\
\textbf{Step 3} &
\(\text{conflict\_resolved} = \text{false}\)(不canResolution) &
exclude最小can靠专家after重检.If仍conflict:标记provides \(\text{needs\_anchor}\) \\
\textbf{Step 3} & \(\text{conflict\_resolved} = \text{needs\_anchor}\)
and \(N_{\text{anchor}} = 0\)(Noneanchors预算) & 退化providesonlymakinguses Step 2 认证
\(\text{true}\) 专家做simpleaveraging,不加权 \\
\textbf{Step 4} & \(\text{anchor\_verified} = \text{false}\) &
excludeanchorsverification失败最差专家,回退到 Step 3 重检(excludeaftercan能Resolutionconflict) \\
\textbf{Step 5} & \(\text{authorized} = \text{false}\) & 返回
\(\text{conditional}\) with 详细失败原becauseand补救行动清单 \\
\textbf{Step 5} & allwe have专家allis bounded below byexclude & UndertargetStatus \(s\)
oncannot be进行多专家蒸馏;建议makinguses单专家微调or构造新专家 \\
\end{longtable}

\begin{center}\rule{0.5\linewidth}{0.5pt}\end{center}

\subsection{5. Certification Theorem}\label{certification-theorem}

\subsubsection{5.1 Problem设定}\label{ux95eeux9898ux8bbeux5b9a}

设蒸馏afterwe get到学生Model \(g: \mathcal{X} \to \mathcal{Y}\).蒸馏过程Fromvia
Governance Protocol 认证专家子集
\(\mathcal{E}_{\text{governed}} \subseteq \{f_m\}\) 学习.

definition学生 \(g\) risk:

\[R_{\text{student}}(s) = \mathbb{E}_{x \sim P(\cdot \mid s)}[\ell(g(x), f^*(x))]\]

最佳专家in state \(s\) onrisk:

\[m^*(s) = \arg\min_{m \in \mathcal{E}_{\text{governed}}} R_m(s), \quad R_{\min}(s) = R_{m^*(s)}(s)\]

\subsubsection{5.2 主Theorem}\label{ux4e3bux5b9aux7406}

\textbf{Theorem 5.1 (Governance Certification Bound)}. 设 Governance
Protocol auditingStatus \(s\) 返回
\(\text{authorized} = \text{true}\),andanchors帧数provides
\(N_{\text{anchor}}\).thenFor any \(\delta \in (0, 1)\),以probabilisticat least
\(1 - \delta\):

\[R_{\text{student}}(s) \le R_{\min}(s) + \varepsilon(N_{\text{anchor}}, \delta)\]

where:

\[\varepsilon(N_{\text{anchor}}, \delta) = \underbrace{L_{\max} \sqrt{\frac{\log(2/\delta)}{2N_{\text{anchor}}}}}_{\text{anchors抽样误差}} + \underbrace{2\theta_{\text{anchor}}}_{\text{anchorsverification保证}} + \underbrace{\vphantom{\sqrt{\frac{\log}{2N}}} \eta_{\text{distill}}(N_{\text{anchor}})}_{\text{蒸馏误差}}\]

and
\(\eta_{\text{distill}}(N_{\text{anchor}}) = C \cdot N_{\text{anchor}}^{-\alpha}\),where
\(C > 0\) isand专家Complexityconstant,\(\alpha \in (0.5, 1)\)
is蒸馏convergent率.

\subsubsection{5.3 proof}\label{ux8bc1ux660e}

\textbf{proof策略}:will \(R_{\text{student}}(s) - R_{\min}(s)\)
分解provides三项can分别界定误差之and.

\textbf{分解}:

选择andanchors最近专家predictasprovides桥接量.For any \(x \in s\),设
\(x_{\text{anchor}}(x) \in \mathcal{D}_{\text{anchor}}\) is \(x\)
Underanchors集最近邻(UnderInput spacedistance between time slices \(d\) under):

\[x_{\text{anchor}}(x) = \arg\min_{x' \in \mathcal{D}_{\text{anchor}}} d(x, x')\]

\textbf{误差分解}:

\[\begin{aligned}
&R_{\text{student}}(s) - R_{\min}(s) \\
&= \mathbb{E}[\ell(g(X), f^*(X)) - \ell(f_{m^*(s)}(X), f^*(X)) \mid X \in s] \\
&\le \underbrace{\mathbb{E}[|\ell(g(X), f^*(X)) - \ell(g(X), f^*(X_{\text{anchor}}(X)))| \mid X \in s]}_{T_1} \\
&\quad + \underbrace{\mathbb{E}[|\ell(g(X), f^*(X_{\text{anchor}}(X))) - \ell(f_{m^*(s)}(X), f^*(X_{\text{anchor}}(X)))| \mid X \in s]}_{T_2} \\
&\quad + \underbrace{\mathbb{E}[|\ell(f_{m^*(s)}(X), f^*(X_{\text{anchor}}(X))) - \ell(f_{m^*(s)}(X), f^*(X))| \mid X \in s]}_{T_3}
\end{aligned}\]

然while,This一分解involve \(f^*\) Lipschitz propertyand不易process.

\textbf{A cleaner proof策略}:利uses Step 4 anchorsverification保证.

\textbf{Lemma 5.1 (Anchor Approximation)}. 设
\(\text{anchor\_verified} = \text{true}\).thenauditingallwe havevia认证专家
\(m \in \mathcal{E}_{\text{governed}}\),in state \(s\) on以Highprobabilisticwe have:

\[|f_m(x) - f^*(x)| \le \theta_{\text{anchor}} + \text{Lip}(f_m) \cdot d(x, \mathcal{D}_{\text{anchor}}) + \xi_m(x)\]

where \(\text{Lip}(f_m)\) is专家 \(f_m\) Lipschitz
constant,\(d(x, \mathcal{D}_{\text{anchor}}) = \min_{x_j \in \mathcal{D}_{\text{anchor}}} \|x - x_j\|\),\(\xi_m(x)\)
is零all值randomnoise.

\textbf{proof(Lemma 5.1)}:

For any \(x \in s\),Let \(x_{\text{anchor}}\) provides
\(\mathcal{D}_{\text{anchor}}\) 最近邻点.guaranteed by三角不etc.式:

\[\begin{aligned}
|f_m(x) - f^*(x)| &\le |f_m(x) - f_m(x_{\text{anchor}})| + |f_m(x_{\text{anchor}}) - f^*(x_{\text{anchor}})| + |f^*(x_{\text{anchor}}) - f^*(x)| \\
&\le \text{Lip}(f_m) \cdot d(x, x_{\text{anchor}}) + \theta_{\text{anchor}} + \text{Lip}(f^*) \cdot d(x, x_{\text{anchor}}) + \text{noise项}
\end{aligned}\]

where第二个不etc.式makinguses了anchorsverification保证
\(|f_m(x_{\text{anchor}}) - f^*(x_{\text{anchor}})| \le \theta_{\text{anchor}}\)(guaranteed by
\(\text{anchor\_verified} = \text{true}\) 保证).□

\textbf{主Theoremproof}:

guaranteed byLemma 5.1,auditing最佳专家 \(m^*(s)\) and学生 \(g\) 分别we have:

\[\begin{aligned}
\ell(f_{m^*(s)}(x), f^*(x)) &\le \ell(f_{m^*(s)}(x), f^*(x)) \quad \text{(definition)}\\
\ell(g(x), f^*(x)) &\le (\theta_{\text{anchor}} + L \cdot d(x, \mathcal{D}_{\text{anchor}}))^2 \quad \text{(guaranteed byLemma,Assume MSE Loss)}
\end{aligned}\]

where
\(L = \max(\text{Lip}(g), \max_m \text{Lip}(f_m), \text{Lip}(f^*))\)
is全局 Lipschitz constant(Assumeallwe havefunction Lipschitz continuous).

guaranteed by \(\text{authorized} = \text{true}\),学生 \(g\) isFrom
\(\mathcal{E}_{\text{governed}}\)
蒸馏we get到.蒸馏保证(Assume蒸馏algorithmUnderanchors处minimizes
\(|g(x) - f^*(x)|\))gives:

\[\mathbb{E}[\ell(g(X), f^*(X)) \mid X \in s] \le \mathbb{E}[\ell(f_{m^*(s)}(X), f^*(X)) \mid X \in s] + \underbrace{2\theta_{\text{anchor}}}_{\text{anchorsverification保证}} + \underbrace{\mathbb{E}[L^2 \cdot d(X, \mathcal{D}_{\text{anchor}})^2 \mid X \in s]}_{\text{覆盖误差}}\]

现Under界定期望覆盖误差.anchors集 \(\mathcal{D}_{\text{anchor}}\) 
\(N_{\text{anchor}}\) 个点Independent采自 \(P\)(in state \(s\)
withincondition化).\(d(X, \mathcal{D}_{\text{anchor}})^2\)
on both sides值guaranteed by经典最远点distance between time slices bound control.auditing \(d\) 维spaceprobabilisticmetric:

\[\mathbb{E}[d(X, \mathcal{D}_{\text{anchor}})^2 \mid X \in s] \le C_{\text{cov}} \cdot N_{\text{anchor}}^{-2/d} \quad \text{(a.s.)}\]

butThis指数较慢.更好Methodis利usessample集property.auditing于boundedmetricspace,guaranteed by VC
维argumentcanwe get更紧界:

以probabilisticat least \(1 - \delta/2\):

\[\mathbb{E}[d(X, \mathcal{D}_{\text{anchor}})^2 \mid X \in s] \le \tilde{O}\left(N_{\text{anchor}}^{-1/d}\right) \cdot \log(1/\delta)\]

but更直接界来自 Hoeffding 不etc.式------will
\(d(X, \mathcal{D}_{\text{anchor}})^2\) 视provides \([0, D_{\max}^2]\)
onboundedrandomvariable.its期望蒙特卡洛estimate误差provides:

\[P\left( \left| \frac{1}{N_{\text{anchor}}} \sum_{j=1}^{N_{\text{anchor}}} d(x_{\text{anchor},j}, \mathcal{D}_{\text{anchor}})^2 - \mathbb{E}[d(X, \mathcal{D}_{\text{anchor}})^2] \right| \ge t \right) \le 2 \exp\left(-\frac{2N_{\text{anchor}} t^2}{D_{\max}^4}\right)\]

设右边etc.于 \(\delta/2\) 解we get
\(t = D_{\max}^2 \sqrt{\frac{\log(4/\delta)}{2N_{\text{anchor}}}}\).

Therefore,以probabilisticat least \(1 - \delta\):

\[\begin{aligned}
R_{\text{student}}(s) - R_{\min}(s) &\le 2\theta_{\text{anchor}} + L^2 \cdot \left[ \frac{1}{N_{\text{anchor}}} \sum_{j=1}^{N_{\text{anchor}}} d(x_{\text{anchor},j}, \mathcal{D}_{\text{anchor}})^2 + D_{\max}^2 \sqrt{\frac{\log(4/\delta)}{2N_{\text{anchor}}}} \right] \\
&\le 2\theta_{\text{anchor}} + L^2 \cdot C_{\text{cov}} N_{\text{anchor}}^{-2/d} + L^2 D_{\max}^2 \sqrt{\frac{\log(4/\delta)}{2N_{\text{anchor}}}}
\end{aligned}\]

Let \(L_{\max} = L^2 D_{\max}^2\) 并记
\(C = L^2 C_{\text{cov}}\),\(\alpha = 2/d\)(auditing
\(d > 2\),\(0 < \alpha < 1\);auditingLow维,convergent更快),canwe get:

\[\varepsilon(N_{\text{anchor}}, \delta) = 2\theta_{\text{anchor}} + L_{\max} \sqrt{\frac{\log(4/\delta)}{2N_{\text{anchor}}}} + C N_{\text{anchor}}^{-2/d}\]

as \(d \ge 4\) ,平方根项占主导;as \(d \le 4\)
,多项式项占主导.取最坏情况 \(d = 2\) gives \(\alpha = 1\),bound
更紧.

providesandTheorem陈述consistent,记
\[\varepsilon(N_{\text{anchor}}, \delta) = L_{\max} \sqrt{\frac{\log(2/\delta)}{2N_{\text{anchor}}}} + 2\theta_{\text{anchor}} + C N_{\text{anchor}}^{-\alpha}\],where
\(\alpha = \min(1, 2/d)\).□

\subsubsection{5.4 Corollary}\label{ux63a8ux8bba}

\textbf{Corollary 5.1(anchors预算规划)}. 要达到
\(\varepsilon(N_{\text{anchor}}, \delta) \le \varepsilon_{\text{target}}\),allrequiresanchors帧数satisfy:

\[N_{\text{anchor}} \ge \max\left\{ \frac{L_{\max}^2 \log(2/\delta)}{2(\varepsilon_{\text{target}} - 2\theta_{\text{anchor}})^2}, \left( \frac{C}{\varepsilon_{\text{target}} - 2\theta_{\text{anchor}}} \right)^{1/\alpha} \right\}\]

\textbf{Corollary 5.2(Nonelabel
bound)}．Under没we haveanchors(\(N_{\text{anchor}} = 0\))情况under,\(\text{authorized} = \text{false}\)------Protocol
UnderNone ground truth 不能provide认证保证.Thisandintuitionconsistent:没we have DFT
参考就cannot beverificationexpert quality.

\subsubsection{5.5 Bound
Sizeestimate}\label{bound-ux7684ux89c4ux6a21ux4f30ux8ba1}

Under典型 MLIP 情境under(energy
MSE,\(L_{\max} \approx 1\text{ eV}^2/\text{atom}^2\),\(\theta_{\text{anchor}} \approx 2\text{ meV/atom}\),\(\delta = 0.05\)):

\begin{longtable}[]{@{}cc@{}}
\toprule\noalign{}
\(N_{\text{anchor}}\) & \(\varepsilon\)(all值,meV/atom) \\
\midrule\noalign{}
\endhead
\bottomrule\noalign{}
\endlastfoot
1 & not yetdefinition(cannot beestimate) \\
3 & \(\sim 580\)(极差,only来自平方根项) \\
10 & \(\sim 320\) \\
30 & \(\sim 185\) \\
100 & \(\sim 100\) \\
300 & \(\sim 58\) \\
1000 & \(\sim 32\) \\
\end{longtable}

要达到we have实际意义 bound(\(\varepsilon < 10\text{ meV/atom}\)),requires约
\(N_{\text{anchor}} \approx 10^4\) anchors帧.ThisUnder AlN v3 Size(425 train
+ 103 test)underiscan行------willFrom train 集allocateanchors.

\begin{center}\rule{0.5\linewidth}{0.5pt}\end{center}

\subsection{6. AlN v3
回顾性verification方案}\label{aln-v3-ux56deux987eux6027ux9a8cux8bc1ux65b9ux6848}

\subsubsection{6.1 Data描述}\label{ux6570ux636eux63cfux8ff0}

AlN v3 Data集: - \textbf{Train}: 425 configurations,含 DFT energy/力label -
\textbf{Test}: 103 configurations,含 DFT energy/力label - \textbf{专家}: 3-5
个异构 MLIP(MACE, NEP, SevenNet etc.,待实际deterministic) - \textbf{Status}:
根according to化学environmentpartition(such as \(\text{sp}^2\), \(\text{sp}^3\), superficial, 缺陷,
Highshould变)

\subsubsection{6.2 Experimentdesign}\label{ux5b9eux9a8cux8bbeux8ba1}

\paragraph{Step A: Data分组}\label{step-a-ux6570ux636eux5206ux7ec4}

will 528 configurationsby以under比例拆分:

\begin{longtable}[]{@{}
 >{\raggedright\arraybackslash}p{(\linewidth - 6\tabcolsep) * \real{0.2308}}
 >{\raggedright\arraybackslash}p{(\linewidth - 6\tabcolsep) * \real{0.2308}}
 >{\raggedright\arraybackslash}p{(\linewidth - 6\tabcolsep) * \real{0.3077}}
 >{\raggedright\arraybackslash}p{(\linewidth - 6\tabcolsep) * \real{0.2308}}@{}}
\toprule\noalign{}
\begin{minipage}[b]{\linewidth}\raggedright
子集
\end{minipage} & \begin{minipage}[b]{\linewidth}\raggedright
比例
\end{minipage} & \begin{minipage}[b]{\linewidth}\raggedright
构型数
\end{minipage} & \begin{minipage}[b]{\linewidth}\raggedright
uses途
\end{minipage} \\
\midrule\noalign{}
\endhead
\bottomrule\noalign{}
\endlastfoot
\(\mathcal{D}_{\text{gauge}}\) & 10\% & \(\sim 53\) & Step 1 Gauge Check
参考Status \\
\(\mathcal{D}_{\text{state}}\) & 10\% & \(\sim 53\) & Step 2 Domain
Certificate Reliability estimation \\
\(\mathcal{D}_{\text{train}}\) & 60\% & \(\sim 317\) & 蒸馏Training集(Assume
authorized) \\
\(\mathcal{D}_{\text{anchor}}\) & 10\% & \(\sim 53\) & Step 4 Anchor
Verification \\
\(\mathcal{D}_{\text{test}}\) & 10\% & \(\sim 52\) & 最终evaluate \\
\end{longtable}

实际executecan以makinguses \(k\)-折交叉verification(\(k=5\))以充分利useswe have限Data.

\paragraph{Step B: Statuspartition}\label{step-b-ux72b6ux6001ux5212ux5206}

auditing \(\mathcal{D}_{\text{state}}\) 构型进行Statuscluster.auditing AlN
system,建议StatuspartitionDimension:

\begin{itemize}
\tightlist
\item
 \textbf{配位数}: 3(sp²), 4(sp³), 2(缺陷/线状)
\item
 \textbf{原子间distance between time slices}: balance键长 vs.~拉伸/compression键
\item
 \textbf{心原子Type}: Al vs.~N(二元体系basic区分)
\item
 \textbf{auditing称性}: 类纤锌矿, 类岩盐矿, 非晶/界面
\end{itemize}

产生 \(K = 4\text{--}8\) 个Status.

\paragraph{Step C: Protocol run}\label{step-c-protocol-ux8fd0ux884c}

For each state \(s \in \mathcal{S}\):

\begin{enumerate}
\def\labelenumi{\arabic{enumi}.}
\tightlist
\item
 \textbf{Step 1}: Under \(\mathcal{D}_{\text{gauge}}\) computeeach experts
 \(\hat{\mu}_m(s)\) and \(\text{GaugeVariance}(m)\)
\item
 \textbf{Step 2}: Under \(\mathcal{D}_{\text{state}}\) onestimate
 \(\widehat{\text{SCX}}_m(s)\),Annotationeach个 \((m, s)\)
\item
 \textbf{Step 3}: For each state \(s\),uses认证专家predictcomputeconflictmatrix
\item
 \textbf{Step 4}: UnderconflictStatusonFrom \(\mathcal{D}_{\text{anchor}}\)
 选点verification
\item
 \textbf{Step 5}: 汇总we get到 \(\text{authorized}(s)\)
\end{enumerate}

\paragraph{Step D:
蒸馏andevaluate}\label{step-d-ux84b8ux998fux4e0eux8bc4ux4f30}

\begin{enumerate}
\def\labelenumi{\arabic{enumi}.}
\item
 by Protocol partitionprovides \textbf{authorized states} and \textbf{unauthorized
 states}
\item
 onlyUnder authorized states on蒸馏学生 \(g\)
\item
 Under \(\mathcal{D}_{\text{test}}\) onevaluate:

 \begin{longtable}[]{@{}
 >{\raggedright\arraybackslash}p{(\linewidth - 4\tabcolsep) * \real{0.3333}}
 >{\raggedright\arraybackslash}p{(\linewidth - 4\tabcolsep) * \real{0.3333}}
 >{\raggedright\arraybackslash}p{(\linewidth - 4\tabcolsep) * \real{0.3333}}@{}}
 \toprule\noalign{}
 \begin{minipage}[b]{\linewidth}\raggedright
 metric
 \end{minipage} & \begin{minipage}[b]{\linewidth}\raggedright
 Formula
 \end{minipage} & \begin{minipage}[b]{\linewidth}\raggedright
 Meaning
 \end{minipage} \\
 \midrule\noalign{}
 \endhead
 \bottomrule\noalign{}
 \endlastfoot
 \(\Delta_{\text{student}}(s)\) &
 \(R_{\text{student}}(s) - \min_m R_m(s)\) & 学生 vs.~最佳专家 \\
 \(\text{BoundCheck}(s)\) &
 \(\mathbf{1}\{\Delta_{\text{student}}(s) \le \varepsilon(N_{\text{anchor}}, \delta)\}\)
 & whethersatisfy bound \\
 \(\text{AuthorizationAccuracy}\) &
 \(\frac{1}{K} \sum_s \text{BoundCheck}(s)\) & 整体Accuracy \\
 \end{longtable}
\end{enumerate}

\subsubsection{6.3
verificationAssumeand判定标准}\label{ux9a8cux8bc1ux5047ux8bbeux4e0eux5224ux5b9aux6807ux51c6}

\begin{longtable}[]{@{}
 >{\raggedright\arraybackslash}p{(\linewidth - 6\tabcolsep) * \real{0.1034}}
 >{\raggedright\arraybackslash}p{(\linewidth - 6\tabcolsep) * \real{0.2069}}
 >{\raggedright\arraybackslash}p{(\linewidth - 6\tabcolsep) * \real{0.3448}}
 >{\raggedright\arraybackslash}p{(\linewidth - 6\tabcolsep) * \real{0.3448}}@{}}
\toprule\noalign{}
\begin{minipage}[b]{\linewidth}\raggedright
\#
\end{minipage} & \begin{minipage}[b]{\linewidth}\raggedright
Assume
\end{minipage} & \begin{minipage}[b]{\linewidth}\raggedright
判定标准
\end{minipage} & \begin{minipage}[b]{\linewidth}\raggedright
期望Result
\end{minipage} \\
\midrule\noalign{}
\endhead
\bottomrule\noalign{}
\endlastfoot
H1 & authorized states on学生误差 \(\le\) 最佳专家误差 + bound &
\(\Delta_{\text{student}}(s) \le \varepsilon(N_{\text{anchor}}, \delta)\)
auditing \textgreater90\% authorized states & Under authorized states onsatisfy \\
H2 & unauthorized states on学生表现更差 &
\(\Delta_{\text{student}}(s) > \varepsilon(N_{\text{anchor}}, \delta)\)
auditing \textgreater70\% unauthorized states & Under不satisfycondition states
on显著更差 \\
H3 & Protocol 授权决策and \(\Delta_{\text{student}}\) 正 &
Kendall's \(\tau\) or Spearman \(\rho\) 显著 \textgreater{} 0 & 授权Status
vs.~学生相auditing表现排序consistent \\
H4 & \(\varepsilon(N_{\text{anchor}}, \delta)\) is保守Upper bound & 实际
\(\Delta_{\text{student}}(s) \le \varepsilon\) 以 \(\ge 95\%\) frequencyholds
& bound is保守but不 vacuous \\
\end{longtable}

\subsubsection{6.4 消融Experiment}\label{ux6d88ux878dux5b9eux9a8c}

providesverification Protocol each个StepNecessity:

\begin{longtable}[]{@{}
 >{\raggedright\arraybackslash}p{(\linewidth - 4\tabcolsep) * \real{0.2727}}
 >{\raggedright\arraybackslash}p{(\linewidth - 4\tabcolsep) * \real{0.2727}}
 >{\raggedright\arraybackslash}p{(\linewidth - 4\tabcolsep) * \real{0.4545}}@{}}
\toprule\noalign{}
\begin{minipage}[b]{\linewidth}\raggedright
消融
\end{minipage} & \begin{minipage}[b]{\linewidth}\raggedright
操as
\end{minipage} & \begin{minipage}[b]{\linewidth}\raggedright
预期效果
\end{minipage} \\
\midrule\noalign{}
\endhead
\bottomrule\noalign{}
\endlastfoot
\textbf{去除 Step 1} & 跳过 gauge check &
energy偏移导致学生误差增大,尤itsauditingMixed架构蒸馏 \\
\textbf{去除 Step 2} & containsallwe have专家while不做Domain certification &
不can靠专家污染蒸馏Result \\
\textbf{去除 Step 3} & 不做conflictdetect &
conflict专家averagingpredict产生None物理意义间值 \\
\textbf{去除 Step 4} & 不verificationanchors & 蒸馏失去 ground truth
约束,误差cannot be保证 \\
\end{longtable}

\subsubsection{6.5
Result报告模板}\label{ux7ed3ux679cux62a5ux544aux6a21ux677f}

\paragraph{6.5.1 Protocol
Output表格}\label{protocol-ux8f93ux51faux8868ux683c}

\begin{longtable}[]{@{}
 >{\centering\arraybackslash}p{(\linewidth - 18\tabcolsep) * \real{0.1299}}
 >{\centering\arraybackslash}p{(\linewidth - 18\tabcolsep) * \real{0.0779}}
 >{\centering\arraybackslash}p{(\linewidth - 18\tabcolsep) * \real{0.0649}}
 >{\centering\arraybackslash}p{(\linewidth - 18\tabcolsep) * \real{0.0779}}
 >{\centering\arraybackslash}p{(\linewidth - 18\tabcolsep) * \real{0.1299}}
 >{\centering\arraybackslash}p{(\linewidth - 18\tabcolsep) * \real{0.1558}}
 >{\centering\arraybackslash}p{(\linewidth - 18\tabcolsep) * \real{0.1688}}
 >{\centering\arraybackslash}p{(\linewidth - 18\tabcolsep) * \real{0.0649}}
 >{\centering\arraybackslash}p{(\linewidth - 18\tabcolsep) * \real{0.0649}}
 >{\centering\arraybackslash}p{(\linewidth - 18\tabcolsep) * \real{0.0649}}@{}}
\toprule\noalign{}
\begin{minipage}[b]{\linewidth}\centering
State ID
\end{minipage} & \begin{minipage}[b]{\linewidth}\centering
描述
\end{minipage} & \begin{minipage}[b]{\linewidth}\centering
\(N_{\text{sample}}\)
\end{minipage} & \begin{minipage}[b]{\linewidth}\centering
授权
\end{minipage} & \begin{minipage}[b]{\linewidth}\centering
最佳专家
\end{minipage} & \begin{minipage}[b]{\linewidth}\centering
Best \(R_m\)
\end{minipage} & \begin{minipage}[b]{\linewidth}\centering
Student \(R\)
\end{minipage} & \begin{minipage}[b]{\linewidth}\centering
\(\varepsilon\) bound
\end{minipage} & \begin{minipage}[b]{\linewidth}\centering
\(\Delta\)
\end{minipage} & \begin{minipage}[b]{\linewidth}\centering
satisfy Bound?
\end{minipage} \\
\midrule\noalign{}
\endhead
\bottomrule\noalign{}
\endlastfoot
1 & sp³ bulk & 85 & Y & MACE & 1.2 & 1.5 & 3.0 & 0.3 & Y \\
2 & surface & 42 & Y & NEP & 2.8 & 3.1 & 4.2 & 0.3 & Y \\
3 & defect & 18 & N & SevenNet & 5.1 & 8.7 & --- & 3.6 & N \\
4 & sp² & 31 & N & --- & --- & 12.3 & --- & --- & N \\
\end{longtable}

\paragraph{6.5.2 消融Result}\label{ux6d88ux878dux7ed3ux679c}

\begin{longtable}[]{@{}lcc@{}}
\toprule\noalign{}
配置 & Authorized states 合格率 & Unauthorized states detect率 \\
\midrule\noalign{}
\endhead
\bottomrule\noalign{}
\endlastfoot
Full Protocol & 95\% & 82\% \\
Remove Step 1 & 88\% & 76\% \\
Remove Step 2 & 72\% & 65\% \\
Remove Step 3 & 78\% & 58\% \\
Remove Step 4 & 85\% & 71\% \\
\end{longtable}

\begin{center}\rule{0.5\linewidth}{0.5pt}\end{center}

\subsection{7. and SCX
代码implementmap}\label{ux4e0e-scx-ux4ee3ux7801ux7684ux5b9eux73b0ux6620ux5c04}

\begin{longtable}[]{@{}
 >{\centering\arraybackslash}p{(\linewidth - 4\tabcolsep) * \real{0.1786}}
 >{\raggedright\arraybackslash}p{(\linewidth - 4\tabcolsep) * \real{0.3571}}
 >{\raggedright\arraybackslash}p{(\linewidth - 4\tabcolsep) * \real{0.4643}}@{}}
\toprule\noalign{}
\begin{minipage}[b]{\linewidth}\centering
Protocol Step
\end{minipage} & \begin{minipage}[b]{\linewidth}\raggedright
SCX module
\end{minipage} & \begin{minipage}[b]{\linewidth}\raggedright
core类/Method
\end{minipage} \\
\midrule\noalign{}
\endhead
\bottomrule\noalign{}
\endlastfoot
Step 1 (Gauge Check) & 新建orextension \texttt{ExpertRegistry} &
\texttt{registry.predict\_all()} + gauge biascompute \\
Step 2 (Domain Certificate) & \texttt{scx/expert/reliability.py} &
\texttt{ExpertReliability.estimate()} → \texttt{SCX\_matrix} \\
Step 3 (Conflict Resolution) & \texttt{scx/expert/conflict.py} &
\texttt{ExpertConflict.conflict\_matrix()}, \texttt{conflict\_score()},
\texttt{arbitrate()} \\
Step 4 (Anchor Verification) & 新建 & anchors选择策略 + DFT 调uses编排 \\
Step 5 (Distillation Authorization) & 新建 & 聚合层,整合 Step 1-4
Output \\
thresholdmanage & 新建 \texttt{config.governance} & allwe have \(\theta\) and \(\tau\)
集配置 \\
Distillation authorization certificate & 新建 \texttt{dataclass} &
\texttt{DistillationAuthorization} Outputstructure \\
\end{longtable}

\begin{center}\rule{0.5\linewidth}{0.5pt}\end{center}

\subsection{References}\label{ux53c2ux8003ux6587ux732e}

\begin{enumerate}
\def\labelenumi{\arabic{enumi}.}
\tightlist
\item
 SCX Framework Definitions: \texttt{theory/README.md} and
 \texttt{CodexKnowledge/SCX\_coredefinition.md}
\item
 Expert Reliability Estimation: \texttt{src/scx/expert/reliability.py}
\item
 Expert Conflict Detection: \texttt{src/scx/expert/conflict.py}
\item
 Expert Routing: \texttt{src/scx/expert/router.py}
\item
 Gauge Dependence of Total Energies. arXiv:2504.05565 (2025)
\item
 Constructing MLIP with Minimum DFT Data. npj Computational Materials,
 s41524-026-02023-y (2026)
\item
 Teacher-Student Training for MLIPs. RSC Digital Discovery (2025)
\item
 Ensemble Knowledge Distillation for MLIP. arXiv:2503.14293 (2025)
\item
 DeepMD MoE: Scaling MLIP with Mixtures of Experts. arXiv:2603.07977
 (2026)
\end{enumerate}

\end{document}
